\section{模型假设}

\begin{enumerate}
  \item 假设附件中的检测记录及字段含义真实可靠，清洗过程中仅对同次采血的重复测序进行合并，不改变原始生物学观测。
  \item 假设同次采血的Y染色体浓度测量误差相互独立、均值为零，且其方差在研究范围内近似稳定，以便利用重复测序差异估计测量误差。
  \item 假设在$10$至$25$孕周的研究区间内，同一BMI水平下的群体达标概率不会随孕周增加而下降；观测中的局部波动主要由测量误差和个体差异造成。
  \item 问题二以首次检测时BMI作为基线BMI，并假设其能够代表孕妇所属的BMI群体，避免使用达标后的BMI造成时间信息倒置。
  \item 假设样本所代表地区的检测流程在研究期内保持一致，模型结论主要用于相似高BMI人群的群体时点建议，不直接替代个体临床诊断。
\end{enumerate}
